% © 2026 Ing. David Jaroš — CC BY-NC-ND 4.0
% Canonical-relation Clifford/Dirac lift, 2026-07-27.
\section{Clifford and generalized-Dirac lift of the canonical relation}
\label{sec:canonical:biquaternion_dirac_lift}

\subsection{Scope: no second tetrad}
The starting point remains the canonical UBT tetrad
\begin{equation}
 E_\mu=\mathcal N_0^{-1/2}D_\mu\Theta\in W_L
\end{equation}
and the central metric identity
\begin{equation}
 \frac12(E_\mu^\sharp E_\nu+E_\nu^\sharp E_\mu)
 =g_{\mu\nu}\mathbf1.
 \label{eq:canonical:metric-before-dirac-lift}
\end{equation}
The construction below introduces no independent vielbein and no spinor-current
replacement for $E_\mu$.  It is an injective operator lift of the same four
biquaternionic covectors.

\subsection{Exact block lift}
Use the Pauli representation
\begin{equation}
 \mathbf e_k=-i_{\rm c}\sigma_k,
 \qquad
 E=i_{\rm c}e^0\mathbf1+e^k\mathbf e_k,
 \qquad e^a\in\mathbb R,
\end{equation}
where $i_{\rm c}$ is the commuting complex unit.  Regard $E$ as a
$2\times2$ complex matrix and define
\begin{equation}
 \boxed{
 \mathcal C(E):=
 \begin{pmatrix}
 0&E\\ E^\sharp&0
 \end{pmatrix}\in\operatorname{Mat}(4,\mathbb C).}
 \label{eq:canonical:block-clifford-lift}
\end{equation}
The map $\mathcal C:W_L\to\operatorname{Mat}(4,\mathbb C)$ is real-linear and
injective.

\paragraph{Theorem (canonical relation implies the curved Clifford relation).}
For all $E,F\in W_L$,
\begin{equation}
 \frac12\{\mathcal C(E),\mathcal C(F)\}
 =H_L(E,F)I_4,
 \label{eq:canonical:block-clifford-identity}
\end{equation}
where
\begin{equation}
 H_L(E,F)\mathbf1
 =\frac12(E^\sharp F+F^\sharp E).
\end{equation}
Consequently, with
\begin{equation}
 \Gamma_\mu:=\mathcal C(E_\mu),
\end{equation}
Eq.~\eqref{eq:canonical:metric-before-dirac-lift} is equivalent to
\begin{equation}
 \boxed{\frac12\{\Gamma_\mu,\Gamma_\nu\}=g_{\mu\nu}I_4.}
 \label{eq:canonical:curved-clifford-relation}
\end{equation}

\begin{proof}
Direct block multiplication gives
\begin{equation}
 \frac12\{\mathcal C(E),\mathcal C(F)\}
 =\frac12
 \begin{pmatrix}
 EF^\sharp+FE^\sharp&0\\
 0&E^\sharp F+F^\sharp E
 \end{pmatrix}.
\end{equation}
For Lorentz paravectors both diagonal expressions equal
$2H_L(E,F)\mathbf1$.  Hence the block matrix is
$H_L(E,F)I_4$.  Injectivity follows because either off-diagonal block recovers
$E$.
\end{proof}


\subsection{Spinor carrier without a new fundamental field}
As a complex vector space, $\mathbb B$ has dimension four.  Fix any complex-linear
isomorphism
\begin{equation}
 \operatorname{vec}:\mathbb B\longrightarrow\mathbb C^4,
 \qquad \Psi:=\operatorname{vec}(\Theta).
\end{equation}
The symbol $\Psi$ therefore denotes the same four-complex-component UBT field
in a column representation; it is not an additional matter field.  The lifted
matrices $\Gamma_\mu=\mathcal C(E_\mu)$ act on this carrier.  A different
choice $\operatorname{vec}'=U\operatorname{vec}$, $U\in GL(4,\mathbb C)$,
gives
\begin{equation}
 \Psi'=U\Psi,
 \qquad
 \Gamma_\mu'=U\Gamma_\mu U^{-1},
\end{equation}
and leaves the Clifford relation, the metric and the characteristic
determinant invariant.  Thus the block representation is a faithful operator
realisation of the canonical biquaternionic relation, not a change of physical
field content.

\subsection{Rank is inherited, not manufactured}
The operator lift adds no geometric degrees of freedom.  The metric is still
that of the tetrad coefficients $e_\mu{}^a$,
\begin{equation}
 g=e\eta e^T,
 \qquad \eta=\operatorname{diag}(-1,1,1,1).
\end{equation}
At every invertible tetrad, the differential
$\delta e\mapsto\delta g$ is surjective onto symmetric $4\times4$ matrices.
Indeed, for an arbitrary symmetric $h$, the choice
\begin{equation}
 \delta e=\frac12h\,g^{-1}e
 \label{eq:canonical:tetrad-right-inverse}
\end{equation}
obeys $\delta g=h$.  The kernel consists precisely of
\begin{equation}
 \delta e=eX,
 \qquad X\eta+\eta X^T=0,
\end{equation}
which is the six-dimensional Lorentz Lie algebra.  Therefore
\begin{equation}
 \boxed{\operatorname{rank}D_e g=10,\qquad\dim\ker D_e g=6.}
\end{equation}
Because $\mathcal C$ is injective, the same rank statement holds when the four
$E_\mu$ are represented by the four $\Gamma_\mu$.

This theorem concerns the map $E\mapsto g$.  The stronger on-shell statement
that the complete UBT equation admits enough solutions of
$E_\mu=D_\mu\Theta/\sqrt{\mathcal N_0}$ remains a separate dynamical theorem.

\subsection{Complex-time / fifth operator channel}
Define the grading matrix
\begin{equation}
 \Gamma_*:=
 \begin{pmatrix}I_2&0\\0&-I_2\end{pmatrix}.
\end{equation}
Every matrix $\mathcal C(E)$ is off diagonal, so
\begin{equation}
 \{\Gamma_*,\mathcal C(E)\}=0,
 \qquad \Gamma_*^2=I_4.
 \label{eq:canonical:grading-anticommutation}
\end{equation}
Thus one may form either
\begin{equation}
 \Gamma_\psi=\Gamma_*,\qquad \Gamma_\psi^2=+I_4,
\end{equation}
or
\begin{equation}
 \Gamma_\psi=i_{\rm c}\Gamma_*,\qquad \Gamma_\psi^2=-I_4,
\end{equation}
according to the signature assigned to the $\psi$ channel.  This establishes
an exact algebraic fifth channel compatible with the four operators derived
from $E_\mu$; it does not by itself decide whether $\psi$ is a physical fifth
coordinate, a compact fibre, or an internal complex-time direction.

The corresponding generalized first-order operator has the candidate form
\begin{equation}
 \boxed{
 \mathcal D_{\rm UBT}[\Theta]\Psi
 =i_{\rm c}\Gamma^\mu[D\Theta]\nabla_\mu\Psi
  +i_{\rm c}\Gamma_\psi D_\psi\Psi
  -\mathcal M[\Theta,D\Theta]\Psi.}
 \label{eq:canonical:generalized-dirac-candidate}
\end{equation}
Equation~\eqref{eq:canonical:generalized-dirac-candidate} is an architectural
candidate until it is derived from the UBT action.  The proved statement is
the Clifford lift and its compatibility with the central metric relation.


\subsection{Principal symbol, causal cone and exact factorisation}
Let
\begin{equation}
 \Gamma^\mu:=g^{\mu\nu}\Gamma_\nu
\end{equation}
and, for a real cotangent vector $\xi=\xi_\mu dx^\mu$, define the
first-order principal symbol
\begin{equation}
 \sigma_4(\xi):=\Gamma^\mu\xi_\mu.
\end{equation}
The curved Clifford relation immediately gives
\begin{equation}
 \boxed{
 \sigma_4(\xi)^2
 =g^{\mu\nu}\xi_\mu\xi_\nu I_4.}
 \label{eq:canonical:principal-symbol-square}
\end{equation}
Since $\sigma_4(\xi)$ is block off diagonal and traceless, its characteristic
polynomial is
\begin{equation}
 \chi_{\sigma_4}(\lambda)
 =\left(\lambda^2-g^{\mu\nu}\xi_\mu\xi_\nu\right)^2,
\end{equation}
and therefore
\begin{equation}
 \boxed{
 \det\sigma_4(\xi)
 =\left(g^{\mu\nu}\xi_\mu\xi_\nu\right)^2.}
 \label{eq:canonical:principal-symbol-determinant}
\end{equation}
Thus the characteristic cone of the canonical generalized-Dirac lift is
exactly the null cone of the metric already defined by the UBT relation.  No
second causal structure is introduced.

Including a fifth anticommuting channel with
$\Gamma_\psi^2=\varepsilon I_4$, $\varepsilon=\pm1$, gives
\begin{equation}
 \sigma_5(\xi,\xi_\psi)
 :=\sigma_4(\xi)+\Gamma_\psi\xi_\psi,
\end{equation}
and the exact factorisation
\begin{equation}
 \boxed{
 \sigma_5(\xi,\xi_\psi)^2
 =\left(g^{\mu\nu}\xi_\mu\xi_\nu
       +\varepsilon\xi_\psi^2\right)I_4.}
 \label{eq:canonical:five-channel-symbol-square}
\end{equation}
Equation~\eqref{eq:canonical:five-channel-symbol-square} is algebraic.  Its
physical interpretation depends on whether $\psi$ is treated as a geometric
coordinate, a compact fibre coordinate, or an internal complex-time spectral
direction.

\subsection{Conditional first-jet rank preservation by the fifth channel}
The fifth matrix is invertible.  If $\Gamma_\psi^2=\varepsilon I_4$, then
\begin{equation}
 \Gamma_\psi^{-1}=\varepsilon\Gamma_\psi.
\end{equation}
Consider a first-order equation written in $\psi$-normal form,
\begin{equation}
 \Gamma_\psi D_\psi\Psi
 +\mathcal F(\Psi,E_0,E_1,E_2,E_3)=0,
 \label{eq:canonical:psi-normal-form}
\end{equation}
where $\mathcal F$ contains no $D_\psi\Psi$.  For every prescribed value of
$\Psi$ and every prescribed Lorentz-admissible four-tetrad $E_\mu$, the unique
pointwise solution is
\begin{equation}
 \boxed{
 D_\psi\Psi
 =-\varepsilon\Gamma_\psi
   \mathcal F(\Psi,E_0,E_1,E_2,E_3).}
 \label{eq:canonical:psi-normal-solution}
\end{equation}
Consequently, at the purely algebraic first-jet level, the equation manifold
of~\eqref{eq:canonical:psi-normal-form} projects surjectively onto the four
spacetime tetrad slots $E_\mu$.  Under these stated hypotheses the equation
therefore does not lower the already proved rank ten of $E\mapsto g$.

This is a conditional pointwise theorem, not yet the full UBT on-shell theorem.
Canonical UBT also uses holomorphy in $\tau=t+i_{\rm c}\psi$.  A strict
Cauchy--Riemann condition relates $D_\psi\Psi$ to the real-time derivative and
removes its independent first-jet status.  The decisive remaining task is to
show either that the covariant complex-time equation possesses a compatible
normal form after holomorphy is imposed, or directly that its constrained jet
projection onto the ten metric directions remains surjective.

\subsection{Remaining proof gates}
A complete closure requires:
\begin{enumerate}
 \item derivation of $\mathcal D_{\rm UBT}$ and $\mathcal M$ from one covariant
       UBT action over $\tau=t+i_{\rm c}\psi$;
 \item proof that the implicit cycle
       $\Theta\to E\to\Gamma\to\omega(E)\to D\Theta$ has local solutions;
 \item preservation of the Lorentz slice and nondegeneracy under evolution;
 \item rank ten after all equations of motion, gauge conditions and
       complex-time constraints are imposed;
 \item derivation of the Einstein low-energy dynamics and the quantum limits.
\end{enumerate}
The historical spinor-current tetrad is not used in these gates.
